\section{Dissertation Overview}
\label{sec:dissertation_overview}

This dissertation investigates the role of structured abstractions in enabling agents to operate effectively under novelty. We begin by developing formal frameworks, taxonomies, and benchmark environments that enable systematic study of novelty as a mismatch between an agent's internal models and the environment. Building on this foundation, we examine how abstractions over actions enable efficient adaptation when execution fails, and how abstractions over interaction enable the learning of skills that generalize across objects, dynamics, and robotic embodiment.

Taken together, these contributions advance a unified perspective: while novelty introduces unavoidable mismatch in open-world environments, the structure of an agent's representations determines whether this mismatch leads to brittle behavior or can be managed through efficient adaptation and generalizable skill learning. By organizing interaction at appropriate levels of abstraction, agents can localize failure, acquire new behaviors when needed, and reuse existing knowledge across diverse conditions.
